
\section{Discussion}
Extending time to a three-dimensional vector subspace allows the Flag-and-Pole framework to recover several foundational structures of relativity and quantum physics from a common temporal geometry. The admissible inertial branch has a pseudo-Euclidean realization with three positive temporal and three negative spatial directions, while the corresponding quadratic displacement norm reduces exactly to the Lorentzian interval. Temporal-frame compatibility then fixes a single causal temporal direction, eliminating the two possible tilt modes and confining temporal translation to the stabilized Pole-time axis. The transverse temporal directions nevertheless remain physically active through rotation within the Plane-of-Now, so the theory retains standard four-dimensional relativistic kinematics and causality without suppressing the underlying three-dimensional structure of time.

This geometric separation yields two scalar measures of evolution, with proper time $\tau$ parametrizing Pole--ground motion and the Flag-time angle $\theta^{t}$ measuring temporal rotation along the normalized direction $\mathbf{i}_{F}=\hat{\vartheta}^{\beta}\mathbf{i}_{\beta}$. Just as the Principle of Relativity requires covariance under changes of the proper-time reference, $\tau\mapsto\tau+\delta\tau$, the same requirement applied to the Flag-time reference, $\theta^{t}\mapsto\theta^{t}+\delta\theta^{t}$, yields gauge covariance. The general field construction accommodates unitary representations, and applying this covariance requirement to the established Abelian and non-Abelian structures naturally introduces the corresponding $U(1)$, $SU(2)$, and $SU(3)$ gauge potentials.

Beyond field representations, the connection between wave and particle descriptions follows from the fundamental relation $dS=\xi\,d\theta^{t}$, which identifies an infinitesimal change of classical action with the corresponding Flag-time rotation through a universal scale. With $P_{\mu}=\partial_{\mu}S$ and $K_{\mu}=\partial_{\mu}\theta^{t}$, this relation gives $P_{\mu}=\xi K_{\mu}$ in the free rest configuration and, by Lorentz covariance, in every inertial frame. The experimentally established Planck--de Broglie relations fix this universal scale as $\xi=\hbar$. Wave--particle complementarity therefore reflects the same free motion expressed through translational action and temporal rotation, with minimal coupling extending this relation to the gauge-covariant temporal-rotor transport law.

From this dynamical foundation, the theory identifies two distinct classes of sources, both of purely geometric origin. First, gauge-covariance geometric sources (GC-sources) arise because the Principle of Relativity requires equations containing derivatives of the temporal rotor to remain covariant under local transformations $g_{F}\in U(V_{F})$, thereby retaining the same form for all inertial observers in the full six-dimensional spacetime. Second, propagation-orthogonal geometric sources (PO-sources) arise algebraically from the rotor action itself rather than from the requirement of covariance of the dynamical equations: the rotor acts trivially on geometric components parallel to transport but remains nontrivial on components orthogonal to $\mathbf{U}$.

These two source mechanisms intersect when the gauge curvature associated with a GC-source possesses a propagation-orthogonal component. Projecting this curvature relative to $\mathbf{U}$ converts it into a PO-source, providing the geometric origin of anomalous-moment interactions. Because this source inherits both the gauge curvature and its propagation-dependent decomposition, its form may be more complex than the forms arising directly from the Plane-of-Now algebra, unless additional symmetry or dynamics ultimately reduces it to a simple closed-form expression.

Finally, although measurement probabilities are not developed here, the six-dimensional geometry suggests a possible basis for them. An observer restricted to the four-dimensional spacetime selected by the causal temporal direction cannot access the complete initial temporal state and can therefore specify it only through statistical inference, making the resulting uncertainty structurally epistemic rather than ontological. Once this initial description is assigned, the rotor, or its wave-function representation, tracks the subsequent evolution of all compatible possibilities. Apart from the hypothesis presented, the theory already yields the positive-definite density $J^{0} = \Psi^{\dagger}\Psi \ge 0$, permitting the usual probability-density interpretation of the wave function to be incorporated into the present framework as an experimental contingency rather than a derived result.
